\renewcommand{\Anno}{2001}

\renewcommand{\NumeroProblema}{1}

\renewcommand{\Testo}{Trovare tutte le coppie $(x,y)$ di numeri reali
tali che 
$$\sqrt[15]{x}-\sqrt[15]{y}=\sqrt[5]{x}-\sqrt[5]{y}=
\sqrt[3]{x}-\sqrt[3]{y}.$$}

\renewcommand{\Risposta}{Le soluzioni sono le infinite coppie del tipo
$(\alpha,\alpha)$, con $\alpha\in\mathbb{R}$, più le seguenti 6 coppie:
$(1,0)$ , $(0,1)$, $(-1,1)$, $(-1,0)$, $(0,-1)$, $(1,-1)$.}

\renewcommand{\Aiutino}{Intanto conviene porre $\sqrt[15]{x}=a$,
$\sqrt[15]{y}=b$, in modo che le equazioni diventino $a-b=a^{3}-b^{3}=a^{5}-b^{5}$.  Scomponiamo i due termini con le terze e quinte potenze.  Se $a\neq b$ un fattore si semplifica ovunque.  A questo punto facendo il quadrato del termine centrale meno il termine
di destra si ottiene \ldots}

\renewcommand{\SoluzioneUfficiale}{
}

\renewcommand{\SoluzioniAlternative}{
}

\renewcommand{\Osservazioni}{
}

\renewcommand{\Commenti}{
}
% Aggiunta redazionale: la soluzione preesistente non viene modificata.
\providecommand{\SoluzioneRedazionale}{}
\renewcommand{\SoluzioneRedazionale}{
% SSSUP-BOZZA-AI: SoluzioneRedazionale; autore=ChatGPT; data=2026-11-08; stato=DA-VERIFICARE
% SSSUP-PROVENIENZA: AI: elaborazione nuova dal testo
Poniamo $a=\sqrt[15]{x}$ e $b=\sqrt[15]{y}$; essendo le radici di indice dispari, questo è lecito per ogni $x,y\in\mathbb R$. Le uguaglianze diventano
\[
a-b=a^3-b^3=a^5-b^5.
\]
Se $a=b$, ogni coppia $(x,y)=(\alpha,\alpha)$ è soluzione. Se $a\neq b$, dividiamo per $a-b$ e otteniamo
\[
a^2+ab+b^2=1,
\qquad a^4+a^3b+a^2b^2+ab^3+b^4=1.
\]
Sottraendo la seconda espressione dal quadrato della prima si trova
\[
ab(a+b)^2=0.
\]
Dunque $a=0$, oppure $b=0$, oppure $a=-b$. Nelle prime due possibilità l'altra variabile vale $1$ o $-1$; nell'ultima vale $a^2=1$. Sono quindi possibili esattamente le sei coppie ordinate di elementi distinti dell'insieme $\{-1,0,1\}$. Elevando alla quindicesima potenza si ottengono le medesime sei coppie per $(x,y)$, e la sostituzione conferma che funzionano tutte.
}
